\subsection{Public Q\&A as a knowledge commons under an AI
shock}\label{ssec:theory_kc_online}\label{ssec:theory_qa_dss}\label{ssec:theory_openinnov}

The IS literature studies why individuals contribute knowledge
to electronic repositories when contribution is voluntary and
costly: motivation is anchored in social capital, reputation,
reciprocity, and self-efficacy rather than direct material
incentives \citep{wasko2005faraj,kankanhalli2005contributing},
and the communities that sustain it require communication flows
above a sustainability threshold or they shrink
\citep{butler2001membership,faraj2011knowledge}. Public
question-and-answer platforms are a specialised
user-generated-content infrastructure in which the unit of
contribution is a problem--solution pair retained by a curation
(acceptance) mechanism as a public, non-rival artefact
\citep{goes2014popularity,temizkankumar2024dss}; reading is
non-rival while writing is reputation-gated, so the production
side is close to a peer-production commons
\citep{benkler2002coase,frischmann2012infrastructure}. Stack
Overflow is the canonical example---the largest public
repository of codified programming know-how
\citep{wachs2022geography}---and is an innovation \emph{input}
(codification, spillover, a contributor pipeline) rather than an
innovation \emph{outcome}; the DDD identifies the former and we
do not claim the latter without independent evidence.

Generative AI is \emph{competing} infrastructure for the same
problem-solving service: a user with a routine bug or how-to
question can now obtain a private answer at near-zero search and
posting cost \citep{leonardi2011affordances,eloundou2024gpts},
and that private conversation produces no public artefact. Two
features make it a candidate disruptor of replenishment.
\emph{Differential answerability}: LLMs handle routine
short-code, how-to, and simple debugging well, but not
version-and-environment configurations or organisation-specific
architecture---the differential the design exploits. \emph{Feedback
to training}: a shrinking flow of human-authored public Q\&A is
also a long-run constraint on the models, since recursive
training on generated content degrades quality
\citep{shumailov2024curse}. Because spillover-based accounts
treat the \emph{type} of codified knowledge as first-order
\citep{cohenlevinthal1990absorptive,jaffe1993geographic}, a
change concentrated on routine substitutable problems is
qualitatively distinct from a uniform change of the same size,
which is why a task-level decomposition is substantively
different from the aggregate decline already documented
\citep{burtch2024chatgpt,delriochanona2024large}.

\subsection{Hypotheses}\label{ssec:hypotheses}\label{ssec:theory_concept}\label{ssec:theory_qa_as_input}\label{ssec:conceptual_significance}

\begin{description}
\item[H1 (substitutability margin).] After the ChatGPT release,
public Q\&A in substitutable question types falls more in tags
whose pre-treatment activity is more LLM-answerable:
$\beta_{\text{DDD}}<0$ in (\ref{eq:ddd}), surviving robustness,
placebo sign-reversal, and pre-trend sensitivity.
\item[H2 (funnel attrition, selection not upgrading).] The shock
propagates through the support funnel with a stable relative
magnitude but a smaller absolute shortfall, and under
fine-grained fixed effects the survivor-composition shift
reflects \emph{selection} rather than within-cell upgrading, so
conditional resolvability shares are null.
\item[H3 (bounded scale).] The DDD estimates only the
LLM-substitutability slice; bounded against an upper-bound
aggregate reference it is small, making visible the order of
magnitude of co-occurring forces the design does not adjudicate.
\end{description}

Downstream propagation (open-source entry, productivity,
innovation outputs) is outside the DDD scope and discussed only
in \S\ref{sec:policy} and \S\ref{sec:limitations} with language
separating supported from speculative claims.
